% GENERATED FILE. Edit canonical lessons, then run npm run generate:books.
% canonical-source-hash: fnv1a64:cfe97f3acf550f51
% canonical-lessons: IT-C170-chitarra, IT-C170-pianoforte, IT-C170-violino, IT-C170-tamburo, IT-C170-pallone

\chapter{Guitar, Piano, Violin, Drum, Ball}
\label{ch:it-guitar-piano-violin-drum-ball}
\hlchaptermodality{170}

\begin{chapteropening}
\textbf{By the end of this chapter:} \emph{I can say a guitar, a piano, a violin, a drum and a ball.}

Complete the last lesson of chapter 170: I can say a guitar, a piano, a violin, a drum and a ball.
\end{chapteropening}

\section[la chitarra]{la chitarra — a guitar}

\label{lesson:IT-C170-chitarra}



\begin{quote}
Before the new one: say the Italian for recycling, then the Italian for energy.
\end{quote}

\subsection*{You'll want to know: la chitarra}
\textbf{la chitarra} — \textquotedblleft{}a guitar\textquotedblright{}.

\begin{cousinweb}
\textbf{chitarra} $\leftarrow$ Arabic \emph{qitara}, \textquotedblleft{}a guitar\textquotedblright{} — borrowed; the Arabic word came from Greek \emph{kithara}, \textquotedblleft{}a lyre\textquotedblright{}. (EtymDB 2.1)
\end{cousinweb}

\begin{grammarlens}[title={What you've built}]
One more word for this chapter.
\end{grammarlens}

\subsection*{Your turn}
\begin{itemize}
\raggedright
  \item \emph{Say it:} \emph{la chitarra}
  \item \emph{Say it:} \emph{la chitarra}, once more
  \item \emph{Say it:} say \emph{la chitarra}
  \item \emph{Recall:} say the Italian for recycling, then the Italian for energy, then say \emph{la chitarra} again
\end{itemize}

\begin{tcolorbox}[breakable,colback=teal!4,colframe=teal!35!black,title={Before you move on}]
What does la chitarra mean? (\textquotedblleft{}A guitar\textquotedblright{}.) Say it once more.
\end{tcolorbox}
\section[il pianoforte]{il pianoforte — a piano}

\label{lesson:IT-C170-pianoforte}



\begin{quote}
Before the new one: say the Italian for energy, then the Italian for a guitar.
\end{quote}

\subsection*{You'll want to know: il pianoforte}
\textbf{il pianoforte} — \textquotedblleft{}a piano\textquotedblright{}.

\begin{cousinweb}
\textbf{pianoforte} is \emph{piano}, \textquotedblleft{}soft\textquotedblright{}, plus \emph{forte}, \textquotedblleft{}loud\textquotedblright{}: the first keyboard that could play both. \emph{forte} $\leftarrow$ Latin \emph{fortis}, \textquotedblleft{}strong\textquotedblright{}. (EtymDB 2.1)
\end{cousinweb}

\begin{grammarlens}[title={What you've built}]
One more word for this chapter.
\end{grammarlens}

\subsection*{Your turn}
\begin{itemize}
\raggedright
  \item \emph{Say it:} \emph{il pianoforte}
  \item \emph{Say it:} \emph{il pianoforte}, once more
  \item \emph{Say it:} say \emph{il pianoforte}
  \item \emph{Recall:} say the Italian for energy, then the Italian for a guitar, then say \emph{il pianoforte} again
\end{itemize}

\begin{tcolorbox}[breakable,colback=teal!4,colframe=teal!35!black,title={Before you move on}]
What does il pianoforte mean? (\textquotedblleft{}A piano\textquotedblright{}.) Say it once more.
\end{tcolorbox}
\section[il violino]{il violino — a violin}

\label{lesson:IT-C170-violino}



\begin{quote}
Before the new one: say the Italian for a guitar, then the Italian for a piano.
\end{quote}

\subsection*{You'll want to know: il violino}
\textbf{il violino} — \textquotedblleft{}a violin\textquotedblright{}.

\begin{cousinweb}
\textbf{violino} — built from \emph{viola}, the stringed instrument, with \emph{-ino}, \textquotedblleft{}little\textquotedblright{}. EtymDB 2.1 derives that \emph{viola} — a different word from the colour — from Old Provençal \emph{viola}, from Medieval Latin \emph{vitula}, \textquotedblleft{}a stringed instrument\textquotedblright{}.
\end{cousinweb}

\begin{grammarlens}[title={What you've built}]
One more word for this chapter.
\end{grammarlens}

\subsection*{Your turn}
\begin{itemize}
\raggedright
  \item \emph{Say it:} \emph{il violino}
  \item \emph{Say it:} \emph{il violino}, once more
  \item \emph{Say it:} say \emph{il violino}
  \item \emph{Recall:} say the Italian for a guitar, then the Italian for a piano, then say \emph{il violino} again
\end{itemize}

\begin{tcolorbox}[breakable,colback=teal!4,colframe=teal!35!black,title={Before you move on}]
What does il violino mean? (\textquotedblleft{}A violin\textquotedblright{}.) Say it once more.
\end{tcolorbox}
\section[il tamburo]{il tamburo — a drum}

\label{lesson:IT-C170-tamburo}



\begin{quote}
Before the new one: say the Italian for a piano, then the Italian for a violin.
\end{quote}

\subsection*{You'll want to know: il tamburo}
\textbf{il tamburo} — \textquotedblleft{}a drum\textquotedblright{}.

\begin{cousinweb}
\textbf{tamburo} $\leftarrow$ Arabic \emph{tunbur} — borrowed. (EtymDB 2.1)
\end{cousinweb}

\begin{grammarlens}[title={What you've built}]
One more word for this chapter.
\end{grammarlens}

\subsection*{Your turn}
\begin{itemize}
\raggedright
  \item \emph{Say it:} \emph{il tamburo}
  \item \emph{Say it:} \emph{il tamburo}, once more
  \item \emph{Say it:} say \emph{il tamburo}
  \item \emph{Recall:} say the Italian for a piano, then the Italian for a violin, then say \emph{il tamburo} again
\end{itemize}

\begin{tcolorbox}[breakable,colback=teal!4,colframe=teal!35!black,title={Before you move on}]
What does il tamburo mean? (\textquotedblleft{}A drum\textquotedblright{}.) Say it once more.
\end{tcolorbox}
\section[il pallone]{il pallone — a ball (football)}

\label{lesson:IT-C170-pallone}



\begin{quote}
Before the new one: say the Italian for a violin, then the Italian for a drum.
\end{quote}

\subsection*{You'll want to know: il pallone}
\textbf{il pallone} — \textquotedblleft{}a ball (football)\textquotedblright{}.

\begin{cousinweb}
\textbf{pallone} — built from \emph{palla}, \textquotedblleft{}a ball\textquotedblright{}, with \emph{-one}, \textquotedblleft{}big\textquotedblright{}. \emph{palla} $\leftarrow$ Lombardic \emph{palla} — borrowed from the Germanic language of the Lombards. (EtymDB 2.1)
\end{cousinweb}

\begin{grammarlens}[title={What you've built}]
One more word for this chapter.
\end{grammarlens}

\subsection*{Your turn}
\begin{itemize}
\raggedright
  \item \emph{Say it:} \emph{il pallone}
  \item \emph{Say it:} \emph{il pallone}, once more
  \item \emph{Say it:} say \emph{il pallone}
  \item \emph{Recall:} say the Italian for a violin, then the Italian for a drum, then say \emph{il pallone} again
\end{itemize}

\begin{tcolorbox}[breakable,colback=teal!4,colframe=teal!35!black,title={Before you move on}]
What does il pallone mean? (\textquotedblleft{}A ball (football)\textquotedblright{}.) Say it once more.
\end{tcolorbox}
